\pdfinfoomitdate=1
\pdftrailerid{}
\pdfsuppressptexinfo=15
\documentclass[11pt,letterpaper]{article}
\usepackage[margin=1in]{geometry}
\usepackage[T1]{fontenc}
\usepackage{lmodern,amsmath,amssymb,amsthm,mathtools,booktabs,microtype,longtable}
\usepackage[colorlinks=true,linkcolor=black,citecolor=black,urlcolor=blue]{hyperref}
\usepackage{enumitem}
\setlist{itemsep=3pt,topsep=5pt}
\setlength{\parskip}{5pt}\setlength{\parindent}{0pt}
\setlength{\emergencystretch}{2em}
\newtheorem{theorem}{Theorem}[section]
\newtheorem{lemma}[theorem]{Lemma}
\newtheorem{proposition}[theorem]{Proposition}
\newtheorem{corollary}[theorem]{Corollary}
\theoremstyle{remark}\newtheorem{remark}[theorem]{Remark}
\newcommand{\Prob}{\mathbb P}\newcommand{\E}{\mathbb E}
\newcommand{\dd}{\,\mathrm d}
\DeclareMathOperator{\NB}{NB}\DeclareMathOperator{\Bin}{Bin}
\DeclareMathOperator{\TV}{TV}
\hypersetup{pdftitle={Report 129 The next logarithmic correction for two inversion sequence classes},pdfauthor={}}
\begin{document}
\begin{center}
{\small REPORT 129}\par\vspace{9pt}
{\LARGE\bfseries The next logarithmic correction}\par
{\LARGE\bfseries for two inversion sequence classes}\par\vspace{8pt}
{\large A279556 and A279551}\par\vspace{8pt}
{\small 2 October 2026}\par
{\small Full quantitative upgrade with a proved input companion}
\end{center}
\begin{abstract}
For the two commitment-tree enumerations A279556 and A279551 we prove
\[
 n\log\mu-\log a(n)=\sigma_D n^{1/3}(\log n)^{2/3}
 \left(1+\frac73\frac{\log\log n}{\log n}
                 +O\!\left(\frac1{\log n}\right)\right),
\]
where $(\mu,D)=(9,1)$ and $(8,1/2)$, respectively, and
$\sigma_D=(3\pi^2/(2D))^{1/3}$. The proof gives a quantitative local
critical root, an anchored derivative estimate, a shrinking-strip
random-clock construction, and an integrable dual penalty. Exact
integer-time extraction and endpoint repair are preserved. We derive
refined inverse thresholds and boundary-derivative growth by explicit
comparison bounds, without differentiating an unknown remainder.
The unchanged Report 127 is included as the proved reference for the
exact combinatorial model, spatial barrier, and uniform discrete
entrance and repair. The reader bundle therefore contains the proofs
of the earlier inputs as well as all new estimates. This is a
logarithmic expansion, not a coefficient equivalent, and it supplies
no effective numerical onset.
\end{abstract}
\setcounter{tocdepth}{1}
{\small\setlength{\parskip}{0pt}
 \let\ReportOriginalAddvspace\addvspace
 \renewcommand{\addvspace}[1]{\ReportOriginalAddvspace{0.3em}}
 \tableofcontents}
\newpage

\section{The theorem and the reader bundle}
An inversion sequence of length $n$ is a sequence $e=(e_1,\ldots,e_n)$
of integers with $0\le e_i<i$. Write $a_{247}(n)$ for the number with
no $i<j<k$ satisfying $e_i\le e_j$ and $e_i\ge e_k$, and $a_{759}(n)$
for the number with no such triple also satisfying $e_j\ne e_k$.
Their forbidden pattern sets are $(000,010,110,120)$ and
$(010,110,120)$, respectively. In the usual indexing,
\[
 a_{247}(n)=\mathrm{A279551}(n),\qquad
 a_{759}(n)=\mathrm{A279556}(n),\qquad a_j(0)=1.
\]
All logarithms are natural. The diffusion parameter $D$ and the
critical exponential growth constant $\mu$ are
\[
\begin{array}{c|ccc}
 j&\mu&D&\sigma_D=(3\pi^2/(2D))^{1/3}\\\hline
 759&9&1&(3\pi^2/2)^{1/3}\\
 247&8&1/2&(3\pi^2)^{1/3}
\end{array}
\]
Set, for sufficiently large $n$,
\begin{equation}\label{eq:scales}
 \begin{gathered}
 N=\log n,\quad J=\log\log n,\quad
 H=n^{2/3}N^{1/3},\quad S=n^{1/3}N^{2/3},\\
 \theta=H/n=S/H,\qquad k_0=S/n=N/H.
 \end{gathered}
\end{equation}
Thus $H^2=nS$, $HS=nN$, and $\theta^2=k_0$.
The symbol $P(n)$ below always denotes a fixed finite power of
$\log n$, with its exponent allowed to increase between estimates.
It never hides a positive power of $n$.

\begin{theorem}\label{thm:main}
For either class, as $n$ tends to infinity through all integers,
\begin{equation}\label{eq:main}
 n\log\mu-\log a_j(n)
 =\sigma_D S+\frac{7\sigma_D}{3}\frac{SJ}{N}
                  +O\!\left(\frac SN\right).
\end{equation}
Equivalently, there are constants $C,n_0$ such that for all $n\ge n_0$,
\begin{equation}\label{eq:envelopes}
 e^{-G(n)-C R(n)}\le a_j(n)\mu^{-n}\le e^{-G(n)+C R(n)},
\end{equation}
where
\[
 \psi(t)=t^{1/3}(\log t)^{2/3},\quad
 R(t)=\frac{\psi(t)}{\log t},\quad
 G(t)=\sigma_D\psi(t)+\frac{7\sigma_D}{3}
                  \frac{\psi(t)\log\log t}{\log t}.
\]
The constants are not asserted to be effective or numerically calibrated.
\end{theorem}
The coefficient $7/3$ comes from adding the $J/6$ contributed by
$\tfrac12\log H$ to the $J$ contributed by $\log\log H$, and
then evaluating the resulting $7J/(6N)$ potential correction along
the singular minimizing curve. Retaining an integrable height error
is necessary to distinguish this coefficient from the remainder.

\subsection{How the earlier input is used}
The bundle contains the unchanged PDF, TeX, checkers, and build source
of Report 127 under \path{earlier-inputs/report127/}. That report
proves the exact counting interpretation and the analytic inputs
stated in Section~\ref{sec:inputs}; its proofs are part of this
reader bundle. It is not necessary to download a previous development
note. The present report gives every new quantitative estimate used
to strengthen the leading logarithmic law.

Statements in the unchanged Report 127 describing the next logarithmic
correction as open refer to the scope of that earlier report. They
are superseded here by Theorem~\ref{thm:main}. Its earlier proof and
files have not been silently revised. The underlying exact commitment
construction is due to Britt and Beaton~\cite{BB}; the arguments here
make no global literature-priority claim.

The finite checks are regression and algebra checks, not replacements for
these proofs. Package and rendering checks have the narrower scopes
recorded in the accompanying README and generated check results.

\section{Exact model and earlier proved inputs}\label{sec:inputs}
\subsection{Counting and normalization}
The two auxiliary label trees start at $(0,0)$ and have edges
\begin{align*}
 (p,c)&\longrightarrow(p+1,c),\\
 (p,c)&\longrightarrow(p+1,c-1)\quad(c>0),\\
 (p,0)&\longrightarrow(p-\ell,b)\quad(0\le\ell\le p-1).
\end{align*}
The last edge has multiplicity
\[
 M_{759}(\ell,b)=C_b\binom\ell b,\qquad
 M_{247}(\ell,b)=C_b\binom{b+1}{\ell-b},\qquad
 C_b=\frac1{b+1}\binom{2b}b.
\]
Out-of-range binomial coefficients are zero. The exact terminal sets are
\[
 \mathcal T_{759}=\{(p,0):p\ge0\},\qquad
 \mathcal T_{247}=\{(p,0):0\le p\le2\}.
\]
Multiplicity-weighted paths ending in these sets count the respective
sequences exactly. In particular the extra auxiliary nodes are not
all valid terminal objects. The completed decorated-history
injectivity proof is in Report 127, Section 2.

Use $(\mu,x_j,y_j)=(9,3,1/2)$ for 759 and $(8,2,1/3)$ for 247,
and normalize an edge of multiplicity $m$ by
\[
 q(s,t)=\frac{m h(t)}{\mu h(s)},\qquad h(p,c)=x_j^p y_j^c.
\]
Every length-$n$ path of multiplicity $m$ has normalized mass
$m\mu^{-n}h(\text{endpoint})$. Thus at a valid terminal its
coefficient multiplier is $x_j^{-p}$. These elementary rows are
substochastic, so the mass at every exact original integer time is
at most one, and $a_j(n)\le\mu^n$ for every $n$.

At a boundary $(p,0)$, a completed commitment with $b\ge1$ has drop
$\ell\le p-1$, followed by $w\ge b$ bulk steps, with its last
fulfillment at the final bulk step. Its total mass, duration, and
displacement are
\begin{equation}\label{eq:cycle}
 Q_j(b,\ell,w)=\mu^{-w-1}x_j^{w-\ell}M_j(\ell,b)
                     \binom{w-1}{b-1},\qquad
 L=w+1,\quad\Delta=w-\ell.
\end{equation}
For $b=0$ use $L=1$, $\Delta=-\ell$, and mass
$\mu^{-1}x_j^{-\ell}M_j(\ell,0)$. The ordinary edge has
$L=\Delta=1$ and mass $a=x_j/\mu$. Every boundary path decomposes
uniquely into these cycles, including the duration-one edges. A
``cycle'' here need not return to the same height.

\subsection{Exact transforms and fixed batch laws}
Put $d_b=C_b/4^b$. Summing unrestricted drops at a fixed $b$ gives
$A_j(\lambda,\eta)d_bR_j(\lambda,\eta)^b$, where
\begin{align}
 A_{759}&=\frac{e^\eta}{9(1-e^{-\lambda}/3)},&
 R_{759}&=\frac{4e^\eta}{9(1-e^{-\lambda}/3)(1-e^{\lambda+\eta}/3)},
 \label{eq:transform759}\\
 A_{247}&=\frac{e^\eta(1+e^{-\lambda}/2)}8,&
 R_{247}&=\frac{e^\eta(1+e^{-\lambda}/2)}{2(1-e^{\lambda+\eta}/4)}.
 \label{eq:transform247}
\end{align}
These follow from geometric and binomial summations in
Report 127, Section 5. They are analytic on a fixed neighborhood of
$(0,0)$, independently of $b$; no unrestricted infinite positive-tilt
sum over $b$ is being asserted.

At a fixed $b$ under the exponential tilt, the drop laws are
\begin{align}
 759:\quad&\ell=b+\NB^{\mathrm{fail}}(b+1,1-e^{-\lambda}/3),\\
 247:\quad&\ell=b+\Bin(b+1,e^{-\lambda}/(2+e^{-\lambda})),
 \label{eq:drop}
\end{align}
where the negative binomial counts failures before the stated number
of successes. Independently, for $b\ge1$, $w$ is a negative-binomial
trials variable with $b$ successes and success probability
$1-e^{\lambda+\eta}/3$ or $1-e^{\lambda+\eta}/4$, respectively.
At $b=0$ duration is one. The legal cutoff conditions only the drop,
so fixed-$b$ independence is preserved; $\Delta$ and $L$ themselves
are not independent.

Let $F_{b,\lambda}$ be the appropriate drop CDF. The exact legal row is
\begin{equation}\label{eq:row}
 T_p(\lambda,\eta)=a e^{\lambda+\eta}
  +A_j(\lambda,\eta)\sum_{b=0}^{p-1}
    d_bR_j(\lambda,\eta)^b F_{b,\lambda}(p-1).
\end{equation}
The sum is empty at $p=0$. The constants are
\[
\begin{array}{c|ccccc}
 j&a&A_0&\alpha&D&r=a+2A_0\\\hline
 759&1/3&1/6&3/2&1&2/3\\
 247&1/4&3/16&4/3&1/2&5/8
\end{array}
\]
In both cases $1-r=2A_0$ and
\begin{equation}\label{eq:logR}
 \log R(\lambda,\eta)=\alpha\eta+\frac{\alpha D}{2}\lambda^2
           +O(|\lambda|^3+|\lambda\eta|+\eta^2).
\end{equation}
For $b\le p$, the tilted drop mean is
$\alpha b+O(p|\lambda|+1)$. Uniformly for $\lambda$ near zero,
\begin{equation}\label{eq:drop-tail}
 \Prob(|\ell-\E\ell|\ge u)
 \le C\exp\{-c\min(u^2/(b+1),u)\}.
\end{equation}
This follows because the centered log moment generating functions of
the geometric or Bernoulli summands are bounded by $C\vartheta^2$
on a fixed interval. At zero tilt, the fixed-$b$ normalized laws also obey
\begin{align}
 \Prob(|\Delta|>z)&\le C e^{-c\min(z^2/(b+1),z)},\label{eq:jump-tail}\\
 \Prob(L>u)&\le C e^{-cu+C(b+1)}.\label{eq:clock-tail}
\end{align}
A Chernoff bound applied to \eqref{eq:logR} proves the first;
a fixed positive duration tilt proves the second. Both include $b=0$.

\subsection{Uniform endpoint and spatial estimates}
We state the discrete inputs with their uniformity explicitly.
Their proofs, including integer-time filling, are in Report 127,
Sections 4, 10, and 12.

\begin{proposition}[Earlier proved inputs]\label{prop:inputs}
There are fixed positive constants with the following properties.
\begin{enumerate}
\item For every sufficiently large integer $p$ there is a legal
normalized path family from the root to $(p,0)$, of deterministic
integer length $E(p)$, satisfying
\begin{equation}\label{eq:entrance-input}
 E(p)\le C p^{3/2}/\sqrt{\log p},\qquad
 \text{mass}\ge e^{-C\sqrt{p\log p}}.
\end{equation}
\item For every sufficiently large integer $p$ and every integer
$m\ge C_0p^{3/2}/\sqrt{\log p}$ there is a legal normalized family
from $(p,0)$ to $(2,0)$ of length exactly $m$, with
\begin{equation}\label{eq:repair-input}
 \text{mass}\ge
 \exp\{-C_1[\sqrt{p\log p}+m\log p/p]\}.
\end{equation}
All elementary multiplicity and bulk-path realizations of its selected
macrocycles are retained.
\item The total normalized mass of all first-hit boundary prefixes
reaching height at least $P$ is at most
\begin{equation}\label{eq:barrier-input}
 C e^{-\kappa\sqrt{(P+2)\log(P+2)}}.
\end{equation}
No restriction on their original duration is needed. The same estimate
bounds the normalized coefficient contribution of length-$n$ paths
whose boundary heights reach $P$.
\end{enumerate}
\end{proposition}
For precision about these references, the entrance and repair proof
uses staircase heights $p_k=\lfloor k^2\log(k+1)\rfloor$ and legal
moderate-deviation cycles of mass at least $e^{-C\log k}$.
Their accumulated time is $O(K^3\log K)$ and logarithmic cost
$O(K\log K)$. Return-cycle blocks and bounded terminal edges fill
every remaining integer duration. Thus the constants do not depend
on a fixed positive lower bound for $p/H$.
The spatial proof uses the increasing concave function
$\sqrt{(p+2)\log(p+2)}$ and a sufficiently small exponential multiple
as a positive supermartingale. First-hit prefix extraction followed
by original-time substochasticity gives the coefficient assertion.
These proof summaries clarify the hypotheses; their full arguments
remain in the supplied unchanged companion.

\section{A quantitative local critical root}
For sufficiently large $p$ and the tilt ranges below, write
$\eta_p(\lambda)$ for the unique solution of
$T_p(\lambda,\eta_p(\lambda))=1$ and $\kappa_p=\eta_p(0)$.
The row is continuous and strictly increasing in $\eta$ near the
brackets constructed below, so those brackets prove existence and
uniqueness of the relevant root.

\begin{lemma}\label{lem:root}
For each fixed $B\ge1$, uniformly for
$|\lambda|\le p^{-1/2}(\log p)^B$, as $p\to\infty$,
\begin{align}
 \kappa_p&=\frac{\tfrac12\log p+\log\log p+c_*
                 +O(\log\log p/\log p)}p,
 &c_*&=\tfrac12\log(\pi/\alpha),\label{eq:scalar}\\
 \eta_p(\lambda)&=\kappa_p-D\lambda^2/2
                         +O(p^{-3/2}(\log p)^{C_B}),\label{eq:anchor}\\
 \eta_p'(\lambda)&=-D\lambda+O(p^{-3/4}(\log p)^{C_B}).
 \label{eq:root-derivative}
\end{align}
For each fixed bounded real $z$, uniformly in this same range,
\begin{equation}\label{eq:shift-row}
 T_p(\lambda,\eta_p(\lambda)-z/p)
              =r+(1-r)e^{-z}+o(1).
\end{equation}
The finite exponent $C_B$ may be enlarged in each appearance.
\end{lemma}
\begin{proof}
First take $\lambda=0$ and set
\[
 s=\tfrac12\log p+\log\log p+c,\quad
 \beta=\log R(0,s/p),\quad
 m=p/\alpha+O(\sqrt p(\log p)^2),
\]
with bounded $c$. Then $\beta=\alpha s/p+O(s^2/p^2)$ and
$\beta m=s+O(p^{-1/2}(\log p)^3)$. By \eqref{eq:drop-tail},
the legal CDF is squeezed between hard cutoffs
$p/\alpha\pm\sqrt p(\log p)^2$, with exponentially small tails.
The endpoint Laplace sum has the quantitative estimate
\begin{equation}\label{eq:laplace}
 \sum_{b=0}^m d_b e^{\beta b}
 =2+\frac{e^{\beta m}}{\sqrt\pi m^{3/2}\beta}
          \left[1+O\!\left(\frac1{\beta m}\right)\right]
            +O(p^{-1/4}(\log p)^C).
\end{equation}
Indeed $d_b\le C(b+1)^{-3/2}$, so the excess from $b\le m/2$
is at most $C\beta\sqrt m e^{\beta m/2}$, of the displayed
polynomially small order. For $b>m/2$, put $b=m-h$ and use
$d_b=(\sqrt\pi b^{3/2})^{-1}(1+O(1/b))$.
The geometric zeroth moment gives the displayed main term, while
its first-moment correction divided by $m$ is
\[
 \frac{C\sum_{h\ge0}h e^{-\beta h}}
      {m\sum_{h\ge0}e^{-\beta h}}=O((m\beta)^{-1}).
\]
The geometric tail, Stirling error, unweighted missing tail
$O(m^{-1/2})$, and integer rounding are smaller.

Near $p/\alpha$, the individual weighted Catalan coefficients are
$O(\log p/p)$. The CDF transition window therefore costs
$O(p^{-1/2}(\log p)^C)$. Substitution in \eqref{eq:row} yields
\begin{equation}\label{eq:scalar-row}
 T_p(0,s/p)=r+2A_0\sqrt{\alpha/\pi}\,e^c
           [1+O(\log\log p/\log p)].
\end{equation}
Since $1-r=2A_0$, monotonicity brackets the root at
$c=c_*+O(\log\log p/\log p)$. A bounded subtraction of $z$ from
$s$ similarly proves \eqref{eq:shift-row} at zero tilt.

For the anchored estimate put $\eta_c=\kappa_p-D\lambda^2/2$.
Analyticity and \eqref{eq:logR} give
\begin{align}
 \log R(\lambda,\eta_c)-\log R(0,\kappa_p)
                         &=O(p^{-3/2}(\log p)^C),\label{eq:Rcomparison}\\
 A(\lambda,\eta_c)-A(0,\kappa_p)
                         &=O(p^{-1/2}(\log p)^C).\nonumber
\end{align}
For example, the unsimplified remainder is bounded by a constant
times $|\lambda|^3+|\lambda|\kappa_p+\kappa_p^2+\lambda^4$.
The cancellation of $D\lambda^2/2$ inside $\log R$ is retained
even when $\eta_c$ is negative.

Choose $h=\sqrt p(\log p)^{B+2}$ and $m_\pm=p/\alpha\pm h$.
The actual mean shift $O(p|\lambda|+1)$ is smaller than $h$.
The two CDFs therefore differ from the hard cutoff only inside this
window, apart from tails bounded by $e^{-c(\log p)^{2B+4}}$.
The entire unrestricted coefficient sum through $p$ is at most a
fixed power of $p$: \eqref{eq:scalar} and \eqref{eq:Rcomparison}
give $\log R=O(\log p/p)$. Thus these tails remain negligible after
multiplication by that sum. The truncated unrestricted sum through
$m_+$ is $O(1)$ by \eqref{eq:laplace}, and each coefficient in the
window is $O(\log p/p)$. The exponent comparison changes any
coefficient with $b\le p$ by relative
$O(p^{-1/2}(\log p)^C)$. Consequently
\begin{equation}\label{eq:rowcomparison}
 |T_p(\lambda,\eta_c)-1|\le Cp^{-1/2}(\log p)^C.
\end{equation}

In an $\eta$ neighborhood of radius $p^{-3/2}(\log p)^{C+1}$,
we have $\partial_\eta T_p\ge cp$. To see this, restrict $b$ to
\begin{equation}\label{eq:active-b}
 [p/\alpha-2p/\log p,\ p/\alpha-p/\log p].
\end{equation}
There are order $p/\log p$ terms with individual legal mass at
least $c\log p/p$. The margin $p/\log p$ dominates both the tilt
mean shift and its fluctuations; changes of the exponent in the
stated neighborhood are $o(1)$. Each of these cycles has
$L\ge b+1\ge cp$. Bracket \eqref{eq:rowcomparison} using this
derivative lower bound to prove \eqref{eq:anchor}. The same comparison
after a bounded shift $-z/p$ proves \eqref{eq:shift-row}.

Finally, $\log T_p$ is convex and strictly increasing in $\eta$.
Its zero sublevel set is the hypograph of the root, so
$\eta_p$ is concave; analyticity and $\partial_\eta T_p>0$ also
make it differentiable. A differentiable concave function has its
derivative between its adjacent secant slopes. Apply these inequalities
to \eqref{eq:anchor} with increment
$h_p=p^{-3/4}(\log p)^C$, taking $C$ sufficiently large and first
enlarging the tilt exponent to $B+1$. Both endpoints then lie in
that enlarged range. The secant error is
\[
 O(h_p)+O(p^{-3/2}(\log p)^C/h_p),
\]
which proves \eqref{eq:root-derivative}. The scalar $\kappa_p$
cancels exactly. Concavity applied only to an unanchored $o(1/p)$
remainder would not supply the needed derivative rate.
\end{proof}

\begin{lemma}[Genuine duration moments]\label{lem:moments}
For the exact row-one probability law
\begin{equation}\label{eq:law}
 \Prob_{p,\lambda}(\text{cycle})
       =Q_p(\text{cycle})e^{\lambda\Delta+\eta_p(\lambda)L},
\end{equation}
in the tilt range of Lemma~\ref{lem:root},
\begin{align}
 cp\le\E L\le Cp,&\qquad \E L^2\le Cp^2,
       &\E\Delta^2&\le Cp(\log p)^C,\label{eq:moments}\\
 \frac{\E\Delta}{\E L}&=D\lambda+O(p^{-3/4}(\log p)^C).
 \label{eq:drift}
\end{align}
\end{lemma}
\begin{proof}
At fixed $b$, the duration law remains negative binomial with success
parameter in a fixed compact subset of $(0,1)$. Its second moment
is $O((b+1)^2)$. Sum against the actual legal $b$-masses, which total
at most one, and use \eqref{eq:active-b} for the lower mean.
For displacement, below $m_+$ the unrestricted fixed-$b$ second
moment is
\[
 O\big((b+1)+(b+1)^2(\lambda^2+\eta_p(\lambda)^2)\big)
 \le Cp(\log p)^C.
\]
Only the unrestricted sum through $m_+$ is used here, and it is
$O(1)$. Above $m_+$, retain legality: $\ell\le p-1$, and $w$ is
independent of that event. The second-moment numerator is bounded
by $Cp^2F_{b,\lambda}(p-1)$, whose exponentially small CDF beats
the polynomial unrestricted mass. The $b=0$ and ordinary terms
have bounded moments. Implicit differentiation of the exact row gives
$\eta_p'(\lambda)=-\E\Delta/\E L$; use
\eqref{eq:root-derivative}. No individual duration is replaced by
its conditional mean and no displacement-duration covariance is dropped.
\end{proof}

\section{The corrected potential and integrable height error}
Define
\begin{equation}\label{eq:An}
 A_n=\frac13+\frac76\frac JN,\qquad e=e_n=N^{-6}.
\end{equation}
On a shrinking strip $ce\le x=p/H\le R$, with fixed $c>0$ and
$R<\infty$, and for $|q|\le P(n)$, Lemma~\ref{lem:root} gives
\begin{align}
 \frac nS\eta_p(\theta q)&=\frac{A_n}{x}-\frac D2q^2+E_n(x,q),
 \label{eq:potential}\\
 |E_n(x,q)|&\le\frac{C(1+|\log x|)}{Nx}+n^{-1/3}P(n).
 \label{eq:potential-error}
\end{align}
Indeed $\lambda\sqrt p=\sqrt{Nx}\,q$, so some fixed $B$ covers
all such tilts. The scalar algebra is
\begin{align*}
 \tfrac12\log p&=\tfrac13N+\tfrac16J+\tfrac12\log x,\\
 \log\log p&=J+\log(2/3)+O((J+|\log x|)/N).
\end{align*}
Divide their sum by $Nx$. The constants $c_*$ and $\log(2/3)$
enter the $1/N$ remainder only; the anchored error in
\eqref{eq:anchor} becomes $n^{-1/3}P(n)$ after multiplication by
$n/S$ on this strip.

\begin{lemma}[Variational profile and height integrability]\label{lem:profile}
For $A$ in a compact neighborhood of $1/3$, the minimum of
\[
 I_A(f)=\int_0^1\left(\frac{f'^2}{2D}+\frac A f\right)\dd t,
 \qquad f(0)=f(1)=0,
\]
over positive absolutely continuous paths of finite action is
\begin{equation}\label{eq:min-value}
 \min I_A=\frac{3A}{M_A}=\sigma_D(3A)^{2/3},\qquad
 M_A=(2DA/\pi^2)^{1/3}.
\end{equation}
The minimizer is parameterized by
\begin{equation}\label{eq:profile}
 f_A(t)=M_A\sin^2\varphi,\qquad
 t=\frac{\varphi-\sin\varphi\cos\varphi}{\pi},
 \quad0\le\varphi\le\pi.
\end{equation}
It satisfies $f_A''=-DA/f_A^2$ and
\begin{equation}\label{eq:profile-integrals}
 \int_0^1\frac{\dd t}{f_A}=\frac2{M_A},\qquad
 \int_0^1\frac{1+|\log f_A|}{f_A}\dd t
 =\frac2{M_A}(1-\log M_A+2\log2)
\end{equation}
when $M_A<1$, as holds here for all sufficiently large $n$.
All these endpoint estimates are uniform for the indicated $A$.
\end{lemma}
\begin{proof}
Since $\dd t=(2/\pi)\sin^2\varphi\dd\varphi$, direct
differentiation proves the ODE and the first integral in
\eqref{eq:profile-integrals}. The ODE first integral is
$f_A'^2/(2D)=A(1/f_A-1/M_A)$; integrating it proves
\eqref{eq:min-value}. The second integral follows from
$\int_0^\pi\log(\sin\varphi)\dd\varphi=-\pi\log2$.
For completeness the latter identity follows by symmetry and the
double-angle formula: if $I=\int_0^{\pi/2}\log\sin u\dd u$, then
$2I=\int_0^{\pi/2}\log(\sin u\cos u)\dd u
=I-(\pi/2)\log2$.

To verify minimality, convexity of $v\mapsto v^2/(2D)$ and of
$x\mapsto A/x$ gives, on $[\varepsilon,1-\varepsilon]$, the
pointwise lower tangent inequality around $f_A$.
The linear terms integrate by parts to a boundary term because
$f_A''=-DA/f_A^2$. For a competing finite-action path $f$, its
zero endpoint and Cauchy--Schwarz imply $f(t)=o(t^{1/2})$ as
$t\downarrow0$ and similarly at one. Since
$f_A(t)\asymp t^{2/3}$ and $f_A'(t)=O(t^{-1/3})$, the boundary
products $f_A'(t)(f(t)-f_A(t))$ vanish. Letting
$\varepsilon\downarrow0$ proves $I_A(f)\ge I_A(f_A)$.
Strict convexity gives uniqueness. These endpoint powers also
show integrability directly and uniformity when $A$ stays in a
compact positive interval.
\end{proof}
A strip supremum bound $|\log x|=O(J)$ would turn
\eqref{eq:potential-error} into an artificial bulk term of size
$J/N$. We instead use the finite integral
\eqref{eq:profile-integrals} in both proof directions.

\section{The coefficient upper bound}
\subsection{A quantitative dual with an integrated penalty}
Take $A=A_n$, $\zeta=N^{-9}$, and $r_n=1-2\zeta$. Set
\[
 u(t)=f_A'(\zeta+r_nt)/D,\qquad
 g(t)=f_A(\zeta+r_nt)/\sqrt{r_n}.
\]
Then $u'=-A/g^2$, $\min g\asymp N^{-6}=e$, and
$\max g<1$ for all sufficiently large $n$. Every fixed derivative
needed below is bounded by $P(n)$. Choose a sufficiently large fixed
$K$ and define $v(0)=0$ by
\begin{equation}\label{eq:dual-v}
 v'(t)=\frac{2A}{g(t)}-\frac D2u(t)^2
       -\frac{K(1+|\log g(t)|)}{N g(t)}.
\end{equation}
The absolute value is smooth because $g<1$. From
Lemma~\ref{lem:profile},
\begin{equation}\label{eq:dual-value}
 v(1)\ge\sigma_D(3A_n)^{2/3}-C/N-C\sqrt e.
\end{equation}
Indeed the unpenalized integral is
\[
 \frac1{r_n}\int_\zeta^{1-\zeta}
 \left(\frac{2A\sqrt{r_n}}{f_A}-\frac{f_A'^2}{2D}\right)\dd t.
\]
Its omitted endpoint pieces are $O(\zeta^{1/3})=O(\sqrt e)$;
the time-rescaling factors cost $O(\zeta)$.
The penalty integrates to $O(1/N)$ by
\eqref{eq:profile-integrals}.

For each $t$ let $F_t(x)=-Ax/g(t)^2+2A/g(t)$. The exact gap is
\begin{equation}\label{eq:dual-gap}
 \frac A x-F_t(x)=\frac A x(1-x/g(t))^2.
\end{equation}
For all $ce\le x\le R$, sufficiently large fixed $K$ makes
\begin{equation}\label{eq:dual-penalty}
 F_t(x)-\frac{K(1+|\log g(t)|)}{N g(t)}
 \le\frac A x-\frac{C(1+|\log x|)}{Nx}-\frac3{Nx},
\end{equation}
where $C$ is from \eqref{eq:potential-error}. If
$1/2\le x/g\le2$, both $1/x$ and $1+|\log x|$ are bounded
by fixed multiples of their counterparts at $g$, so the penalty
alone suffices. Outside that range \eqref{eq:dual-gap} is at
least $A/(4x)$, which absorbs the displayed error since
$|\log x|=O(J)$ and $J/N\to0$. The polynomially small part of
\eqref{eq:potential-error} is also less than $1/(Nx)$ uniformly.
Consequently
\begin{equation}\label{eq:dual-root}
 u'(t)x+v'(t)\le\frac nS\eta_p(\theta u(t))-\frac2{Nx}.
\end{equation}
After multiplying by $k_0=S/n$, the slack is exactly $2/p$.

\subsection{Uniform rows including very small heights}
For integer original time $0\le t\le n$ and $0\le p\le RH$, put
\[
 \Phi(t,p)=S[u(t/n)p/H+v(t/n)].
\]
Retain only legal cycles with $t+L\le n$ and endpoints in this
height region, multiplying their masses by
$e^{\Phi(t+L,p+\Delta)-\Phi(t,p)}$. Call a cycle good if
$L\le d_n n$ and $|\Delta|\le d_n H$, where $d_n=n^{-1/12}$.
Taylor's formula gives
\begin{equation}\label{eq:taylor}
 \Phi(t+L,p+\Delta)-\Phi(t,p)
 \le\theta u(t/n)\Delta+
 k_0[u'(t/n)p/H+v'(t/n)+d_nP(n)]L.
\end{equation}
To include the mixed term explicitly, after dividing by $S$ and
writing $l=L/n$, $d=\Delta/H$, $x=p/H$, its remainder is
\begin{align*}
 &[u(t/n+l)-u(t/n)]d\\
 &\quad+x[u(t/n+l)-u(t/n)-u'(t/n)l]\\
 &\quad+[v(t/n+l)-v(t/n)-v'(t/n)l],
\end{align*}
bounded by $P(n)(l|d|+l^2)$. This proves \eqref{eq:taylor}.

At $p\ge ceH$, its extra $\eta$ shift times $p$ is at most
$RN d_nP(n)=o(1)$. Equation~\eqref{eq:dual-root} leaves at least
$1/p$ root slack. The shifted-row limit \eqref{eq:shift-row} with
$z=1$ bounds every such good-cycle row by a fixed number below one.

Choose the fixed constant $c>0$ sufficiently small for the remaining
heights. At $p\le ceH$, the effective tilts in \eqref{eq:taylor}
obey
\[
 \lambda^2\le Ck_0/e,\qquad |\eta|\le Ck_0/e,\qquad
 \beta=\log R(\lambda,\eta)\le Ck_0/e,
\]
with constants independent of small $c\le1$. Indeed
$|u'|p/H\le C/e$ there, and the terms of $v'$ satisfy the
same bound. If $\beta\ge0$, discarding only the CDF gives
\begin{equation}\label{eq:low-excess}
 \sum_{b=0}^{p}d_b(e^{\beta b}-1)
 \le C\beta\sqrt{p+1}e^{\beta p}
 \le n^{-1/3+Cc}P(n).
\end{equation}
This uses $e^z-1\le ze^z$ and
$\sum_{b\le p}bd_b\le C\sqrt{p+1}$. For negative $\beta$
there is no excess. Take $c$ so that $Cc<1/6$.
The prefactor and ordinary-edge perturbations vanish uniformly,
so these rows are bounded above by $r+o(1)$. At $p=0$ there is
only the ordinary edge. Equality with $r+o(1)$ is not claimed at
fixed finite height, and no local-root approximation is used there.

For bad cycles, $|\Phi|\le SP(n)$ on the retained region.
The zero-tilt fixed-$b$ estimates
\eqref{eq:jump-tail}--\eqref{eq:clock-tail}, retaining $b\le RH$
and using $\sum d_b\le2$, give transformed mass at most
\begin{equation}\label{eq:bad}
 Ce^{SP(n)}\left\{
 e^{-c\min(d_n^2H^2/(RH+1),d_nH)}
 +e^{-cd_nn+C(RH+1)}\right\}=o(1).
\end{equation}
The first negative exponent has order $n^{1/2}$ times a fixed
logarithmic power, and the second has order $n^{11/12}$ minus
$O(H)$. Both dominate $SP(n)$. Thus all transformed rows are
bounded by a single fixed $\rho<1$ eventually.

\subsection{Exact original time and large heights}
The transformed mass over every possible number of cycles is at
most $1/(1-\rho)$. Keeping only paths with actual final original
time $n$ can only decrease it. Telescoping the transform and then
applying the genuine terminal multiplier gives
\[
 e^{-\Phi(n,p)}x_j^{-p}
 =e^{-Sv(1)}e^{-p(\log x_j+\theta u(1))}\le e^{-Sv(1)},
\]
since $\theta P(n)\to0$. Hence the confined normalized
coefficient is at most $Ce^{-Sv(1)}$.

For paths with a boundary height beyond $RH$, the spatial input
\eqref{eq:barrier-input} gives contribution
\[
 \exp\{-(\kappa\sqrt{2R/3}+o(1))S\}.
\]
Fix $R$ sufficiently large that the coefficient exceeds
$\sigma_D+1$. This is negligible even at the new logarithmic
scale. Finally,
\[
 \sigma_D(3A_n)^{2/3}
 =\sigma_D+\frac{7\sigma_D}{3}\frac JN+O(J^2/N^2),
 \qquad J^2/N^2=O(1/N).
\]
Together with \eqref{eq:dual-value}, $\sqrt e=N^{-3}$, and the
fixed geometric prefactor, this proves
\begin{equation}\label{eq:upper-coeff}
 n\log\mu-\log a_j(n)
 \ge\sigma_D S+\frac{7\sigma_D}{3}\frac{SJ}{N}-O(S/N).
\end{equation}

\section{The coefficient lower bound with a shrinking strip}
\subsection{Entrance and the proposed central path}
Take $P_n=\lfloor eH\rfloor$ and apply
\eqref{eq:entrance-input}. Write $a_n=E(P_n)/n$. Since
$\log(eH)/N\to2/3$,
\begin{equation}\label{eq:entrance-scale}
 a_n\le Ce^{3/2},\qquad
 \text{entrance logarithmic cost}\le C\sqrt e\,S.
\end{equation}
Let $f=f_{1/3}$ and choose the left endpoint time $t_e$ with
$f(t_e)=e$. For a sufficiently large fixed $A_{\rm rep}$, put
\begin{equation}\label{eq:path}
 \delta=A_{\rm rep}e^{3/2},\quad b_n=1-\delta,\quad
 c_n=\frac{1-2t_e}{b_n-a_n},\quad
 f_n(t)=f(t_e+c_n(t-a_n)).
\end{equation}
Increase $A_{\rm rep}$ so that $a_n\le\delta/2$ and the
uniform repair threshold below holds. Because $t_e=O(e^{3/2})$,
$c_n=1+O(e^{3/2})$. On $[a_n,b_n]$ the path has both endpoints
$e$, lies in $[e,M_{1/3}]$, and has every required fixed derivative
bounded by $P(n)$. Extend it beyond $b_n$ by the same formula for
time $n^{-1/12}$; this is valid eventually because the distance
to its singular endpoint has only a negative logarithmic order.

\subsection{Controlled law and the clock}
At an active departure time $s$, use \eqref{eq:law} with
\begin{equation}\label{eq:control}
 \lambda_s=\theta q(s/n),\qquad q=f_n'/D.
\end{equation}
Stop at the first boundary time $s_K\ge nb_n$, on leaving
$[eH/2,RH]$, or after
\begin{equation}\label{eq:cap}
 m_n=\left\lceil\frac{4n}{c eH}\right\rceil+1
\end{equation}
cycles, where $c$ is the lower-mean constant in
\eqref{eq:moments}. All increments below are set to zero after
stopping. The chosen $R$ is larger than the profile's maximum.
All active departure states satisfy the quantitative local lemmas.

The clock martingale $Z_k=\sum_{i<k}(L_i-\E[L_i\mid\text{past}])$
satisfies
\[
 \E Z_{m_n}^2\le m_n C H^2\le nH P(n).
\]
Doob's inequality therefore yields
\begin{equation}\label{eq:clock-control}
 \max_k|Z_k|/n\le n^{-1/12}
\end{equation}
with failure probability at most $n^{-1/6}P(n)$. Also
\begin{equation}\label{eq:mesh}
 \E\sum_i(L_i/n)^2\le n^{-1/3}P(n),\qquad
 \sum_i(L_i/n)^2\le n^{-1/6}
\end{equation}
with failure probability at most $n^{-1/6}P(n)$. On this event
the maximum time mesh is at most $n^{-1/12}$.
The cycle cap cannot be the first stop without the desired time
crossing or spatial exit: absent a prior stop,
$\E[L_i\mid\text{past}]\ge ceH/2$, and the cap in
\eqref{eq:cap} gives accumulated conditional mean above $2n$.
Its actual elapsed time would still be below $n$, contradicting
\eqref{eq:clock-control}. The factor $4$ accounts for the
lower strip height $eH/2$.

Put $v_i=\E[\Delta_i\mid\text{past}]/\E[L_i\mid\text{past}]$.
The anchored derivative estimate gives, uniformly,
\begin{equation}\label{eq:physical-drift}
 v_i=\theta f_n'(s_i/n)+O(n^{-1/2}P(n)),\qquad
 v_i/\theta=f_n'(s_i/n)+O(n^{-1/6}P(n)).
\end{equation}
The reward martingale $M_k=\sum_{i<k}(\Delta_i-v_iL_i)$ obeys
$\E M_{m_n}^2\le nP(n)$. To check the rate without independence,
use $(\Delta-vL)^2\le2\Delta^2+2v^2L^2$ in
\eqref{eq:moments}; each active increment has second moment at
most $HP(n)$ and $m_n\le(n/H)P(n)$.
Doob's inequality gives
\begin{equation}\label{eq:reward-control}
 \max_k|M_k|/H\le n^{-1/12}
\end{equation}
with failure probability at most $n^{-1/6}P(n)$.

Since the total time is at most $n+\max_iL_i$, these estimates
imply at every visited boundary
\begin{equation}\label{eq:path-sum}
 p_k/H=p_0/H+\sum_{i<k} f_n'(s_i/n)L_i/n
                     +O(n^{-1/12}P(n)).
\end{equation}
The Riemann error in this sum is at most
$P(n)\sum_i(L_i/n)^2$. The entrance rounding is $O(1/H)$.
It follows, on an event whose probability tends to one, that
\begin{equation}\label{eq:tube}
 \sup_k|p_k/H-f_n(s_k/n)|\le\tau_n:=n^{-1/14},\qquad
 0\le s_K/n-b_n\le n^{-1/12}.
\end{equation}
Here $n^{-1/12}P(n)=o(n^{-1/14})$. Since $\tau_n=o(e)$,
the tube excludes a spatial exit: the proposed curve has margin
at least $e/2$ from the lower boundary and fixed upper margin.
The stop is therefore the intended original-time crossing.
This proves a quantitative shrinking-strip result rather than
appealing to a fixed-strip convergence theorem.

\subsection{Exact likelihood and an integrated error}
For each realized stopped prefix its inverse change of measure is
exact:
\begin{equation}\label{eq:likelihood}
 \prod_i Q_{p_i}(\text{cycle}_i)
 =\prod_i\Prob_{p_i,\lambda_{s_i}}(\text{cycle}_i)e^{-B_n},
 \qquad
 B_n=\sum_i[\lambda_{s_i}\Delta_i+
                    \eta_{p_i}(\lambda_{s_i})L_i].
\end{equation}
No duration is rounded or frozen. Equations
\eqref{eq:potential}--\eqref{eq:potential-error}, evaluated at the
actual random heights on \eqref{eq:tube}, give
\begin{align}\label{eq:likelihood-expand}
 \frac{B_n}S={}&\sum_iq(s_i/n)\frac{p_{i+1}-p_i}{H}\nonumber\\
 &+\sum_i\left(\frac{A_n}{p_i/H}-\frac D2q(s_i/n)^2\right)
                   \frac{L_i}n+O(1/N)+n^{-c_*'}P(n)
\end{align}
for some positive $c_*'$. The $O(1/N)$ is obtained by integration,
not by a strip supremum. Specifically,
\begin{equation}\label{eq:integrated-remainder}
 \sum_i\frac{1+|\log(p_i/H)|}{p_i/H}\frac{L_i}n=O(1).
\end{equation}
On $[e/2,R]$ this height function is Lipschitz, with constant
$P(n)$ (also across $x=1$). Its path-tube and mesh errors are
polynomially small; its integral along $f_n$ is uniformly bounded
by \eqref{eq:profile-integrals} and the time change in
\eqref{eq:path}. The short overshoot contributes only
$n^{-1/12}P(n)$. This proves \eqref{eq:integrated-remainder}.
The anchored-root remainder contributes $n^{-1/3}P(n)$.

For the first sum in \eqref{eq:likelihood-expand}, discrete
summation by parts bounds the replacement error by
\[
 \tau_n\big(|q(\text{start})|+|q(\text{end})|+\TV(q)\big)
 \le\tau_nP(n).
\]
Thus no bounded-total-variation assertion for the random boundary
path is needed. Riemann errors for the deterministic smooth path
are polynomially small. The second sum is treated by its tube,
mesh, and lower-height bounds. Hence
\begin{equation}\label{eq:likelihood-action}
 \frac{B_n}S=\int_{a_n}^{b_n}
      \left(\frac{f_n'^2}{2D}+\frac{A_n}{f_n}\right)\dd t
                  +O(1/N)+n^{-c_*'}P(n).
\end{equation}
The overshoot again has time at most $n^{-1/12}$ and every relevant
integrand there is at most $P(n)$.
The controlled event has probability tending to one, and is eventually
at least $1/2$. Integrating \eqref{eq:likelihood} over it costs
only an additional factor $1/2$ in the lower mass bound.

\subsection{Every integer time and evaluation of the action}
On the tube, for sufficiently large $n$,
\[
 eH/2\le p_K\le2eH,\qquad
 \delta n/2\le n-s_K\le\delta n.
\]
Choose the fixed $A_{\rm rep}$ large enough that the repair input
\eqref{eq:repair-input} applies throughout this state and time range.
It appends a legal family of exactly $n-s_K$ steps ending at $(2,0)$,
with logarithmic cost at most
\begin{equation}\label{eq:repair-cost}
 C\big(\sqrt e\,S+(\delta/e)S\big)\le C\sqrt e\,S.
\end{equation}
The entrance cost has the same order. The terminal multiplier is
exactly the constant $x_j^{-2}$. Concatenation is injective by the
completed-history argument in Report 127: the entrance duration is
deterministic, the interior stop is the first boundary visit at or
after $nb_n$, and the repair macrocycle pattern is a deterministic
function of its start state and residual time. Distinct decorated
prefixes remain distinct completed histories. Every internal
multiplicity realization counted by $Q$ is retained.

Set
\[
 K_e=\int_{t_e}^{1-t_e}\frac{f'^2}{2D}\dd t,
 \qquad W_e=\int_{t_e}^{1-t_e}\frac{\dd t}{f}.
\]
The central action in \eqref{eq:likelihood-action} is exactly
$c_nK_e+A_nc_n^{-1}W_e$. The endpoint estimates give
\[
 K_e+\tfrac13W_e=\sigma_D+O(\sqrt e),\quad
 W_e=2\sigma_D+O(\sqrt e),\quad c_n=1+O(e^{3/2}).
\]
Therefore
\begin{equation}\label{eq:lower-action}
 \int_{a_n}^{b_n}\left(\frac{f_n'^2}{2D}+\frac{A_n}{f_n}\right)\dd t
 =\sigma_D+\frac{7\sigma_D}{3}\frac JN+O(\sqrt e).
\end{equation}
The unperturbed profile suffices because the action is linear in
$A_n$ at that path. Combine \eqref{eq:entrance-scale},
\eqref{eq:likelihood-action}, \eqref{eq:repair-cost}, and
\eqref{eq:lower-action}. Since $\sqrt e=N^{-3}=O(1/N)$ and
every $n^{-c}P(n)=o(1/N)$, this proves
\begin{equation}\label{eq:lower-coeff}
 n\log\mu-\log a_j(n)
 \le\sigma_D S+\frac{7\sigma_D}{3}\frac{SJ}{N}+O(S/N).
\end{equation}
Together with \eqref{eq:upper-coeff}, this proves
Theorem~\ref{thm:main} at every sufficiently large integer $n$.

\section{Error ledger and the scope of the new term}
The following ledger lists the errors that must stay below $SJ/N$.
Unless indicated otherwise, the stated errors are relative to $S$.
\begin{longtable}{@{}p{0.26\textwidth}p{0.27\textwidth}p{0.41\textwidth}@{}}
\toprule
Source & Bound & Reason it is harmless \\\midrule
\endfirsthead
\toprule
Source & Bound & Reason it is harmless \\\midrule
\endhead
Scalar local root & $O(1/N)$ after integration & The height error is $(1+|\log x|)/(Nx)$ and is integrated along the profile \\
Anchored tilt remainder & $n^{-1/3}P(n)$ & Uniform on the strip down to $eH$ \\
Controlled drift & $n^{-1/6}P(n)$ & Obtained from anchored concavity, not differentiation of an unspecified error \\
Clock fluctuation & variance $\le nHP(n)$ & Threshold $n^{11/12}$ has failure probability $n^{-1/6}P(n)$ \\
Reward fluctuation & variance $\le nP(n)$ & Threshold $Hn^{-1/12}$ has failure probability $n^{-1/6}P(n)$ \\
Squared time mesh & expectation $n^{-1/3}P(n)$ & Threshold $n^{-1/6}$ controls Riemann errors and overshoot \\
Path tube & radius $n^{-1/14}$ & Smaller than every fixed negative logarithmic power and therefore than the strip margin \\
Upper Taylor remainder & $n^{-1/12}P(n)$ per unit time & Its root shift times $p$ is $o(1)$, preserving $1/p$ slack \\
Very small heights & row excess $n^{-1/3+Cc}P(n)$ & Choose fixed $c$ with $Cc<1/6$; the row is bounded above by $r+o(1)$ \\
Entrance and repair & $O(N^{-3})$ cost; $O(nN^{-9})$ time & Uniform discrete constructions retain every residual integer duration \\
Likelihood height error & $O(1/N)$ & Lipschitz tube and mesh comparison with an integrable deterministic profile \\
Large jumps and durations & negative exponents of orders $n^{1/2}$ and $n^{11/12}$ & Both dominate $SP(n)$ even for the growing dual \\
Boundary heights above $RH$ & exponent greater than $(\sigma_D+1)S$ & One fixed sufficiently large $R$ suffices \\
Terminal and counting factors & fixed multiplier and $O(1)$ logarithmic cost & Exact terminal $(2,0)$ and injective concatenation \\
Expansion of minimum & $O(J^2/N^2)$ & This is $O(1/N)$ \\\bottomrule
\end{longtable}
The estimate leaves an unknown coefficient at scale
$S/N=n^{1/3}(\log n)^{-1/3}$. Exponentiating such a remainder does
not give a multiplicative $1+o(1)$ error. In particular neither an
exact multiplier nor a polynomial prefactor is determined. The
presence of $c_*$ in the scalar root does not isolate the next global
coefficient: other accumulated effects at the same scale have not
been evaluated. No ratio theorem, full transseries, effective onset,
or finite-data calibration is asserted.

\section{Refined inverse thresholds}
In this and the next section write $\sigma=\sigma_D$ and
$\lambda=\log\mu$. Let $\psi$, $R$, and $G$ be as in
Theorem~\ref{thm:main}.
\begin{corollary}\label{cor:inverse}
For real $T\to\infty$, define
$N(T)=\min\{n\ge0:a_j(n)\ge T\}$ and $s=(\log T)/\lambda$.
Then
\begin{equation}\label{eq:inverse}
 N(T)=s+\frac\sigma\lambda\psi(s)
       +\frac{7\sigma}{3\lambda}
          \frac{\psi(s)\log\log s}{\log s}
       +O\!\left(\frac{\psi(s)}{\log s}\right).
\end{equation}
This conclusion requires neither monotonicity of the coefficients
nor a ratio asymptotic.
\end{corollary}
\begin{proof}
Put $m(s)=s+G(s)/\lambda$ and consider integers within
$O(\psi(s))$ of $s$. Direct differentiation of the explicit smooth
functions gives $G'(t)=O(\psi(t)/t)$ and
$R'(t)=O(R(t)/t)$. Uniformly in this range,
\[
 G(n)=G(s)+O(\psi(s)^2/s)=G(s)+o(R(s)),\qquad
 R(n)=R(s)(1+o(1)).
\]
The small-error assertion follows from
$\psi(s)\log s/s\to0$. For a fixed sufficiently large $K$, let
$n_+=\lceil m(s)+KR(s)\rceil$. The coefficient lower envelope
in \eqref{eq:envelopes} implies
$\log a_j(n_+)\ge\lambda s$ eventually. Integer rounding costs
$O(1)=o(R(s))$.

For the other direction, $a_j(n)\le\mu^n$ excludes every
integer $n<s$. The explicit function
$t\mapsto\lambda t-G(t)+CR(t)$ is strictly increasing eventually,
since its derivative tends to $\lambda>0$. At
$m(s)-KR(s)$ it is less than $\lambda s$ for large enough
fixed $K$. The upper envelope therefore excludes every integer
between $s$ and this threshold. These two brackets prove
$N(T)=m(s)+O(R(s))$, which is \eqref{eq:inverse}.
\end{proof}

\section{Refined boundary derivative growth}
Let $b_n=a_j(n)\mu^{-n}$, $F(z)=\sum_{n\ge0}a_j(n)z^n$,
$\rho=1/\mu$, and $f(t)=F(\rho e^{it})$. Theorem~\ref{thm:main}
gives radius $\rho$ and $\sum n^k b_n<\infty$ for every fixed
integer $k$. Thus $f$ is smooth, all its derivatives are represented
by uniformly convergent series, and positivity gives
\begin{equation}\label{eq:moment-norm}
 \|f^{(k)}\|_\infty=M_k:=\sum_{n\ge0}n^k b_n\quad(k\ge1).
\end{equation}
Its maximum is attained at $t=0$. Radial boundary derivatives are
$F^{(k)}(\rho)=\mu^k\sum_{n\ge k}(n)_k b_n$, with the falling
factorial $(n)_k=n(n-1)\cdots(n-k+1)$.

\begin{corollary}\label{cor:moments}
As $k\to\infty$ through integers,
\begin{align}
 \log\|f^{(k)}\|_\infty
 ={}&3k\log k-2k\log\log k
       +[\log(2D/\pi^2)-3]k\nonumber\\
 &\hspace{8mm}-\frac{k\log\log k}{\log k}
                     +O\!\left(\frac{k}{\log k}\right),
 \label{eq:moment-refinement}\\
 \log F^{(k)}(\rho)
 ={}&3k\log k-2k\log\log k
       +[\log\mu+\log(2D/\pi^2)-3]k\nonumber\\
 &\hspace{8mm}-\frac{k\log\log k}{\log k}
                     +O\!\left(\frac{k}{\log k}\right).
 \label{eq:radial-refinement}
\end{align}
These are consequences of the explicit envelopes
\eqref{eq:envelopes}; they assert no extra coefficient term.
\end{corollary}
\begin{proof}
We first optimize explicit smooth bounds. For any fixed real $c$,
put $V_c(x)=G(x)+cR(x)$ for sufficiently large real $x$. With
$y=\log x$, $a=7/3$, and $b=c/\sigma$,
\begin{equation}\label{eq:explicit-potential}
 V_c(e^y)=\sigma e^{y/3}y^{2/3}
             \left(1+\frac{a\log y+b}{y}\right).
\end{equation}
Its first and second $y$ derivatives are eventually positive,
with second derivative asymptotic to
$(\sigma/9)e^{y/3}y^{2/3}$. Thus $ky-V_c(e^y)$ is strictly
concave eventually, and its unique eventual maximum tends to
infinity. Finitely many smaller $x$ have objective only $O(k)$.
We differentiate only this explicit potential.

Write $B(y)=a\log y+b$. The stationary equation is
\begin{equation}\label{eq:stationarity}
 k=\frac\sigma3e^{y/3}y^{2/3}
       \left(1+\frac{B(y)+2}{y}
                    +O\!\left(\frac{\log y}{y^2}\right)\right).
\end{equation}
Taking logarithms and multiplying by three gives, with
$L=\log k$,
\begin{equation}\label{eq:saddle-equation}
 y+2\log y+\frac{7\log y+3b+6}{y}
 =3L+3\log(3/\sigma)
                +O\!\left(\frac{(\log y)^2}{y^2}\right).
\end{equation}
First this shows $y\sim3L$; then it shows
$y=3L-2\log L+O(1)$. Substituting this back and expanding yields
\begin{equation}\label{eq:saddle-location}
 y=3L-2\ell+c_0-\frac\ell L+O(1/L),\qquad
 \ell=\log L,\quad c_0=\log(3/\sigma^3).
\end{equation}
Here is the coefficient calculation explicitly. If
$y=3L-2\ell+c_0+d\ell/L+O(1/L)$, then $2\log y$ contributes
$-4\ell/(3L)$ at this order, and $7\log y/y$ contributes
$7\ell/(3L)$. Their combined $\ell/L$ coefficient in
\eqref{eq:saddle-equation} is $d+1$, so $d=-1$.
The remaining residual is $O(1/L)$ and the left derivative tends
to one, justifying the stated remainder by the mean-value theorem.

Divide \eqref{eq:explicit-potential} by
\eqref{eq:stationarity}. At the saddle,
\begin{equation}\label{eq:saddle-value-potential}
 V_c(e^y)=3k\left(1-\frac2y
                    +O\!\left(\frac{\log y}{y^2}\right)\right)
          =3k-\frac{2k}{L}+O(k\ell/L^2).
\end{equation}
Consequently, for every fixed $c$,
\begin{align}\label{eq:optimized}
 \sup_{x\ge x_0}\{k\log x-V_c(x)\}
 ={}&3kL-2k\ell+(c_0-3)k-\frac{k\ell}{L}+O(k/L).
\end{align}
The saddle has
$x\asymp k^3/(\log k)^2$. Rounding to a nearest integer changes
the objective by $o(1)$: in a fixed neighborhood of that saddle,
the absolute $x$ derivative is $O(k/x)$, and $k/x\to0$.
Thus \eqref{eq:optimized} also holds for the integer supremum.

Now use \eqref{eq:envelopes}. A single saddle integer for
$V_C$ supplies the lower bound for $M_k$.
For the upper bound reserve the summable positive factor $e^{-R(n)}$:
\[
 n^k b_n\le e^{-R(n)}
       \exp\{k\log n-G(n)+(C+1)R(n)\}
       \qquad(n\ge n_0).
\]
The series $\sum_{n\ge n_0}e^{-R(n)}$ is finite, and the remaining
supremum is \eqref{eq:optimized} with $c=-(C+1)$.
The earlier coefficients contribute $e^{O(k)}$, negligible beside
the saddle on the logarithmic scale $k/L$. Therefore
\eqref{eq:optimized} proves \eqref{eq:moment-refinement}, since
$3/\sigma_D^3=2D/\pi^2$.

For radial derivatives, $(n)_k\le n^k$ gives the upper bound.
At the same lower-bound saddle $n\asymp k^3/(\log k)^2$,
\[
 \log\frac{(n)_k}{n^k}
 =\sum_{j=0}^{k-1}\log(1-j/n)=O(k^2/n)=o(k/\log k).
\]
The falling factorial therefore preserves the displayed refinement;
the factor $\mu^k$ adds $k\log\mu$.
\end{proof}
The unknown coefficient of $R(n)$ affects the moment only at
$O(k/\log k)$, which explains why the additional term in
\eqref{eq:moment-refinement} is nevertheless determined. In the
convention $\|f^{(k)}\|_\infty\le C A^k(k!)^s$, the sharp
Gevrey order remains $3$, as already established in Report 127.

\section{Verification and reproducibility}
The report supplies every new quantitative estimate, and the unchanged
earlier companion supplies the complete proofs of the stated inputs.
A mathematical argument, a finite regression check, a byte-integrity
check, and a rendered-page review have distinct scopes; none is a
substitute for another.

The reader bundle is local and offline. It contains this report's
PDF and editable TeX, the unchanged Report 127 package, and finite-check and reproduction tools.
The outer closed inventory rejects missing or changed files, extra
files or directories, unsafe inventory paths, and symbolic links.
The earlier package retains its own original inventory.
The README gives the precise checker commands, their scope, and
the build environment. Full replay verifies checks in normal and
optimized Python, rebuilds the report in clean external directories,
extracts a fresh deterministic ZIP, repeats the checks, and compares
the repack. All runtime products stay outside the immutable package.

The finite companion tests explicit identities and bounded examples.
It does not certify asymptotic uniformity, the singular variational
argument, martingale convergence, or infinite tail estimates.
A source inventory is a record relative to its supplied verifier,
not an authenticity signature against simultaneous malicious
replacement. Different TeX or compression versions may produce
equivalent content with different bytes; the recorded environment
is used for the bytewise replay.

\begin{thebibliography}{9}
\bibitem{R127}
\emph{Report 127. Sharp logarithmic deficit constants for two inversion
sequence classes}, 2 October 2026. Unchanged proof and source companion
included under \path{earlier-inputs/report127/}. Exact counting:
Sections 2--3. Exact transforms and concentration: Section 5.
Uniform entrance and every-integer-time repair: Section 10.
Spatial barrier: Section 12.
\bibitem{BB}
Nathan Britt and Nicholas Beaton,
\emph{Completing the enumeration of inversion sequences avoiding triples
of relations}, arXiv:2512.21943v3, 29 September 2026.
Exact commitment trees: Sections 3.8--3.9, equations (3.40)--(3.45),
Lemmas 5--6 and Appendix A.
\url{https://arxiv.org/html/2512.21943v3}.
\bibitem{OEIS247}
OEIS Foundation Inc., \emph{A279551}.
\url{https://oeis.org/A279551}. Sequence identifier.
\bibitem{OEIS759}
OEIS Foundation Inc., \emph{A279556}.
\url{https://oeis.org/A279556}. Sequence identifier.
\end{thebibliography}
\end{document}
